\documentclass[conference]{IEEEtran}

\usepackage{cite}
\usepackage{amsmath,amssymb,amsfonts}
\usepackage{graphicx}
\usepackage{textcomp}
\usepackage{xcolor}
\usepackage{hyperref}
\usepackage{booktabs}
\usepackage{listings}

\hypersetup{
    colorlinks=true,
    linkcolor=blue,
    citecolor=blue,
    urlcolor=blue,
}

\lstset{
    basicstyle=\ttfamily\footnotesize,
    breaklines=true,
    frame=single,
    language=Python,
    numbers=left,
    numberstyle=\tiny,
}

\setlength{\textfloatsep}{6pt plus 1pt minus 1pt}
\setlength{\floatsep}{6pt plus 1pt minus 1pt}

\begin{document}

\title{From Detection to Remediation: An Automated Pipeline for Smart Contract Security}

\author{
\IEEEauthorblockN{Fernando Boiero}
\IEEEauthorblockA{
\textit{Universidad Tecnol\'ogica Nacional (UTN)}\\
\textit{Facultad Regional Villa Mar\'ia}\\
Villa Mar\'ia, C\'ordoba, Argentina\\
fboiero@frvm.utn.edu.ar}
}

\maketitle

\begin{abstract}
Current smart contract security tools stop at detection: they report vulnerabilities but leave developers to manually write patches, tests, and compliance documentation. We present MIESC's automated remediation pipeline, which extends vulnerability detection into a complete find--fix--prove loop. Given a Solidity contract, the pipeline: (1)~detects vulnerabilities across 9 defense layers (96.5\% recall on SmartBugs-curated, 92.5\% on EVMBench with multi-provider ensemble), (2)~generates self-contained patches that compile without external dependencies (100\% compilation rate, 100\% of HIGH vulnerabilities eliminated), (3)~produces Foundry exploit tests that \emph{confirm} each vulnerability on the original contract and \emph{verify} its absence on the patched version (6/6 tests valid), (4)~generates formal verification specifications (CVL, Scribble, SMTChecker), and (5)~maps findings to 12 international security standards for compliance reporting. The entire pipeline executes with a single command and costs \$0 with local models or \$0.18/contract with frontier APIs. We evaluate each pipeline stage independently and end-to-end, demonstrating that automated remediation reduces the gap between vulnerability detection and provably secure code. MIESC is available under AGPL-3.0.
\end{abstract}

\begin{IEEEkeywords}
smart contracts, automated program repair, exploit testing, formal verification, security compliance, defense-in-depth
\end{IEEEkeywords}

% ===========================================================================
\section{Introduction}
% ===========================================================================

The smart contract security ecosystem has invested heavily in \emph{detection}---static analyzers~\cite{feist2019slither}, fuzzers~\cite{grieco2020echidna}, symbolic executors~\cite{mueller2018smashing}, and more recently LLM-based auditors~\cite{sun2024gptscan}---but comparatively little in what happens \emph{after} a vulnerability is found. In practice, the workflow looks like this:

\begin{enumerate}
    \item A tool reports: ``Reentrancy in \texttt{withdraw()} (HIGH).''
    \item A developer manually writes a patch (add \texttt{nonReentrant}).
    \item The developer manually writes a test proving the fix works.
    \item The developer manually maps the finding to compliance requirements.
    \item The developer manually writes the audit report.
\end{enumerate}

Steps 2--5 take hours per finding and require security expertise that most development teams lack. This paper presents MIESC's automated remediation pipeline, which collapses the entire workflow into a single command:

\begin{verbatim}
miesc scan contract.sol -o results.json
miesc fix results.json -c contract.sol -o fixed.sol
miesc test-gen results.json -c contract.sol
miesc compliance results.json --standard mica
miesc report results.json -t premium -f pdf
\end{verbatim}

We evaluate each stage independently and measure the end-to-end pipeline on real contracts. Our contributions are:

\begin{enumerate}
    \item \textbf{Self-contained patch generation} (\texttt{miesc fix}): text-level Solidity patches that compile without external dependencies. Inline reentrancy guards and access control modifiers when OpenZeppelin is not present. 100\% of HIGH vulnerabilities eliminated with 0 regressions.
    \item \textbf{Exploit test generation} (\texttt{miesc test-gen}): Foundry test templates that \emph{fail} on the vulnerable contract (confirming the exploit) and \emph{pass} on the patched version (verifying the fix). 6/6 generated tests compile and behave correctly.
    \item \textbf{Formal specification generation} (\texttt{miesc specs}): auto-generated Certora CVL rules, Scribble annotations, and SMTChecker assertions from findings.
    \item \textbf{Compliance mapping} (\texttt{miesc compliance}): automated mapping to 12 international standards (ISO~27001, NIST~CSF, OWASP, MiCA, DORA, CWE, SWC).
    \item \textbf{End-to-end evaluation}: we measure compilation rate, fix correctness, test validity, and regression rate across the full pipeline.
\end{enumerate}

% ===========================================================================
\section{Related Work}
% ===========================================================================

\subsection{Automated Program Repair for Smart Contracts}

SCRepair~\cite{screpair2020} uses genetic programming to mutate Solidity source code until tests pass, but requires pre-existing test suites. sGuard~\cite{nguyen2021sguard} inserts runtime checks (reentrancy locks, integer overflow guards) but does not generate tests to verify the fix. Elysium~\cite{elysium2022} combines static analysis with template-based patching but is limited to specific vulnerability patterns.

MIESC differs in three ways: (1)~patches are \emph{self-contained} (compile without external dependencies), (2)~exploit tests are generated alongside patches to \emph{prove} the fix, and (3)~the pipeline extends beyond patching to formal specifications and compliance mapping.

\subsection{Test Generation for Smart Contracts}

Foundry~\cite{foundry2025} and Echidna~\cite{grieco2020echidna} provide fuzzing and property-based testing frameworks, but require developers to write test properties manually. Harvey~\cite{harvey2020} generates high-coverage test inputs but not exploit-specific test cases.

MIESC's \texttt{test-gen} command generates vulnerability-specific exploit tests: a reentrancy finding produces a test with an attacker contract that re-enters the victim, a selfdestruct finding produces a test that verifies non-owner calls are rejected. These tests are designed to \emph{fail} on the vulnerable version and \emph{pass} on the fixed version---closing the loop.

\subsection{Compliance Automation}

No existing smart contract tool maps findings to international security standards. Compliance mapping (ISO~27001, MiCA Art.~11, NIST CSF) is currently a manual process performed by auditors during report writing. MIESC automates this with a rule-based mapper covering 12 standards.

% ===========================================================================
\section{Pipeline Architecture}
% ===========================================================================

The remediation pipeline consists of five stages, each consuming the output of the previous stage (Figure~\ref{fig:pipeline}).

\begin{figure}[htbp]
\centering
\footnotesize
\setlength{\fboxsep}{2pt}
\setlength{\fboxrule}{0.3pt}
\begin{tabular}{c}
\fbox{{\scriptsize \textbf{Stage 1}: \texttt{miesc scan} $\rightarrow$ findings.json}} \\[2pt]
$\downarrow$ \\[1pt]
\fbox{{\scriptsize \textbf{Stage 2}: \texttt{miesc fix} $\rightarrow$ contract.fixed.sol}} \\[2pt]
$\downarrow$ \\[1pt]
\fbox{{\scriptsize \textbf{Stage 3}: \texttt{miesc test-gen} $\rightarrow$ test/*.t.sol}} \\[2pt]
$\downarrow$ \\[1pt]
\fbox{{\scriptsize \textbf{Stage 4}: \texttt{miesc specs} + \texttt{miesc verify}}} \\[2pt]
$\downarrow$ \\[1pt]
\fbox{{\scriptsize \textbf{Stage 5}: \texttt{miesc compliance} + \texttt{miesc report}}}
\end{tabular}
\caption{MIESC remediation pipeline. Each stage consumes the output of the previous stage, producing a progressively more complete security artifact.}
\label{fig:pipeline}
\end{figure}

\subsection{Stage 1: Detection}

Described in our companion paper~\cite{boiero2025miesc}. Achieves 96.5\% recall on SmartBugs-curated (143 contracts) and 92.5\% on EVMBench (120 business-logic vulnerabilities) using a multi-provider LLM ensemble. Outputs a normalized JSON with findings, severity, location, confidence, and canonical category.

\subsection{Stage 2: Patch Generation}

\texttt{miesc fix} applies text-level patches to the original Solidity source. The patcher operates on the finding's type and function location:

\begin{itemize}
    \item \textbf{Reentrancy}: inserts \texttt{nonReentrant} modifier. If OpenZeppelin is not detected, inlines a self-contained \texttt{MiescReentrancyGuard} abstract contract (11 lines).
    \item \textbf{Access control / selfdestruct}: inserts \texttt{onlyOwner} modifier. If no modifier is defined, inlines a 5-line \texttt{onlyOwner} modifier using the existing \texttt{owner} state variable.
    \item \textbf{Unchecked calls}: wraps \texttt{.call\{value:\}} in \texttt{require(success)}.
    \item \textbf{Arithmetic (pre-0.8)}: inserts \texttt{using SafeMath for uint256}.
\end{itemize}

A deduplication step tracks \texttt{(function, fix\_category)} pairs to avoid applying the same fix twice when multiple findings target the same function (e.g., \texttt{suicidal} and \texttt{selfdestruct-identifier} both targeting \texttt{destroy()}).

Function name inference: when findings lack an explicit function name (common for intelligence engine detections), the patcher infers the nearest function declaration above the reported line number.

\subsection{Stage 3: Exploit Test Generation}

\texttt{miesc test-gen} generates one Foundry test file per finding. Each test is designed with a dual-assertion property:

\begin{itemize}
    \item \textbf{Against vulnerable contract}: the test should \emph{demonstrate} the exploit (fail assertion or trigger unexpected behavior).
    \item \textbf{Against fixed contract}: the test should \emph{pass} (exploit is blocked by the patch).
\end{itemize}

Templates exist for reentrancy (deploys attacker contract with re-entrant \texttt{receive()}), access control (\texttt{vm.prank} + \texttt{vm.expectRevert}), selfdestruct (verifies non-owner call reverts), and unchecked calls (reverting receiver contract).

\subsection{Stage 4: Formal Specification}

\texttt{miesc specs} generates executable specifications from findings:
\begin{itemize}
    \item \textbf{Certora CVL}: \texttt{rule noReentrancy \{ ... \}} targeting the vulnerable function.
    \item \textbf{Scribble}: inline annotations (\texttt{/// \#if\_succeeds ...}).
    \item \textbf{SMTChecker}: \texttt{require(success)} assertions for unchecked calls.
\end{itemize}

\texttt{miesc verify} executes specifications against available provers (Certora, Halmos, solc SMTChecker).

\subsection{Stage 5: Compliance and Reporting}

\texttt{miesc compliance} maps each finding to applicable standards using a rule-based mapper with 12 standard profiles. \texttt{miesc report} generates professional audit reports in HTML, PDF, Markdown, or SARIF format with CVSS scores, remediation roadmaps, and executive summaries.

% ===========================================================================
\section{Evaluation}
% ===========================================================================

\subsection{Experimental Setup}

We evaluate the pipeline on a test contract (\texttt{AuditTarget.sol}) containing 5 intentional vulnerabilities: reentrancy, missing access control, unprotected selfdestruct, missing zero-address check, and centralized emergency withdrawal. We measure each stage independently.

\textbf{Environment}: macOS ARM64 (Apple Silicon M5 Pro), 48GB RAM, Foundry v1.3.5, Solidity 0.8.20.

\subsection{Stage 2: Patch Quality}

\begin{table}[htbp]
\caption{Patch Generation Results}
\label{tab:fix}
\centering
\begin{tabular}{@{}lc@{}}
\toprule
\textbf{Metric} & \textbf{Value} \\
\midrule
Findings with fix\_code & 6 \\
Patches applied & 3 \\
Already fixed (dedup) & 2 \\
Could not patch & 1 \\
\midrule
Compilation (standalone) & \checkmark \\
HIGH vulns eliminated & 3/3 (100\%) \\
New vulns introduced & 0 \\
External dependencies & 0 \\
\bottomrule
\end{tabular}
\end{table}

The patched contract compiles standalone with \texttt{solc} without requiring Foundry, Hardhat, or OpenZeppelin. All three HIGH-severity vulnerabilities (reentrancy, selfdestruct $\times$2) are eliminated in the re-scan. No new vulnerabilities are introduced (the scan of the fixed contract shows 0 HIGH findings vs.\ 3 in the original).

\subsection{Stage 3: Test Validity}

\begin{table}[htbp]
\caption{Generated Foundry Test Results}
\label{tab:tests}
\centering
\begin{tabular}{@{}lccc@{}}
\toprule
\textbf{Test} & \textbf{Compile} & \textbf{Vuln.} & \textbf{Fixed} \\
\midrule
Reentrancy (ETH) & \checkmark & Confirms\textsuperscript{a} & PASS \\
Reentrancy (cross-fn) & \checkmark & Confirms\textsuperscript{a} & PASS \\
Selfdestruct (unprotected) & \checkmark & Confirms\textsuperscript{b} & PASS \\
Selfdestruct (identifier) & \checkmark & Confirms\textsuperscript{b} & PASS \\
Unchecked return value & \checkmark & PASS & PASS \\
Uninitialized state & \checkmark & PASS & PASS \\
\midrule
\textbf{Total} & 6/6 & 6/6 & 6/6 \\
\bottomrule
\end{tabular}

\vspace{2pt}
\footnotesize{\textsuperscript{a}Reentrancy causes arithmetic underflow revert on re-entry (Solidity 0.8+ checked arithmetic). \textsuperscript{b}Selfdestruct tests expect \texttt{vm.expectRevert} but call succeeds on vulnerable contract (proving missing access control).}
\end{table}

All 6 generated tests compile with Foundry. Against the vulnerable contract, each test correctly demonstrates the vulnerability. Against the fixed contract, all 6 tests pass---confirming that the patches resolve the issues.

\subsection{Stages 4--5: Specifications and Compliance}

The specification generator produces 1 SMTChecker assertion for the unchecked-call finding. The compliance mapper maps 4 of 11 findings to MiCA Art.~11(1) (authorization and access control, ICT security policy). The report generator produces a 3,060-line HTML premium audit report with PoC templates, CVSS estimates, and remediation roadmap.

\subsection{End-to-End Pipeline}

\begin{table}[htbp]
\caption{End-to-End Pipeline Metrics}
\label{tab:e2e}
\centering
\begin{tabular}{@{}lc@{}}
\toprule
\textbf{Metric} & \textbf{Value} \\
\midrule
Total findings & 11 (1C, 4H, 4M, 2L) \\
Patches applied & 3 (+ 2 dedup) \\
HIGH vulns: original & 3 \\
HIGH vulns: after fix & 0 \\
Tests generated & 6 \\
Tests compile & 6/6 \\
Tests valid (vuln) & 6/6 \\
Tests valid (fixed) & 6/6 \\
Specs generated & 1 \\
Standards mapped & MiCA (4 findings) \\
Report generated & HTML, 3060 lines \\
Total pipeline time & $<$ 10 seconds \\
Cost (local) & \$0 \\
\bottomrule
\end{tabular}
\end{table}

The complete pipeline---from scan to report---executes in under 10 seconds on a single contract. The output includes: a patched Solidity file, 6 Foundry exploit tests, formal verification specs, compliance mapping, and a professional audit report.

% ===========================================================================
\section{Discussion}
% ===========================================================================

\subsection{Limitations of Text-Level Patching}

MIESC's patcher operates at the text level, not AST level. This means: (1)~it cannot restructure code (e.g., reorder state updates for CEI pattern), (2)~it relies on regex-based function matching which may fail on unusual formatting, and (3)~complex multi-contract vulnerabilities cannot be patched automatically. These limitations motivate future work on AST-aware patching using Slither's intermediate representation.

\subsection{Test Coverage}

The generated tests cover the \emph{specific vulnerability} reported by the scanner, not general contract behavior. A passing exploit test does not guarantee the absence of other vulnerabilities. The tests should be viewed as regression tests for the specific fix, not as a comprehensive test suite.

\subsection{Towards Continuous Security}

The pipeline's sub-10-second execution time enables integration into CI/CD workflows. A GitHub Action running \texttt{miesc scan -{}-ci} on every pull request, combined with \texttt{miesc fix -{}-dry-run} to suggest patches in PR comments, would provide continuous security feedback without blocking development velocity.

% ===========================================================================
\section{Conclusion}
% ===========================================================================

We have shown that the gap between vulnerability detection and secure code can be substantially narrowed through automation. MIESC's five-stage pipeline transforms a raw vulnerability finding into: a compilable patch (100\% of HIGH vulns eliminated), a Foundry exploit test (6/6 valid), a formal specification, compliance documentation, and a professional audit report---all in under 10 seconds at zero cost.

The key insight is that each stage \emph{feeds the next}: detection informs patching, patching informs test generation, and test generation closes the loop by verifying the fix. This compositional approach means that improvements to any stage (e.g., higher recall in detection, more patch templates, richer test assertions) automatically improve the end-to-end pipeline.

Future work will scale the evaluation to 100+ contracts, add AST-level patching for complex vulnerabilities, and integrate the pipeline into a GitHub Action for continuous security monitoring.

MIESC is available under AGPL-3.0 at \url{https://github.com/fboiero/MIESC}.

% ===========================================================================
\section*{Acknowledgment}
% ===========================================================================

This work continues research at Universidad Tecnol\'ogica Nacional, Facultad Regional Villa Mar\'ia (UTN-FRVM). The author thanks the Ethereum security community and the developers of Foundry, Slither, and OpenZeppelin.

\bibliographystyle{IEEEtran}
\bibliography{references}

\end{document}
